\section{Live Demos 与行动边界}

课堂用订机票、移动端操作、驾驶测试和 Action API 展示 agent 的可见行动能力。高质量阅读不能只复述“它做到了”，而要逐步问：目标如何输入、环境返回了什么、何处发生状态改变、哪些步骤需要确认、失败时怎样停止或恢复。

\subsection{订机票 Demo：从请求到 itinerary}

第一组五张画面记录一次 browser-agent flight-booking 流程。起始页先给出任务语境，随后开发界面发出请求，浏览器进入搜索结果、checkout，最后显示 itinerary。连续状态说明 agent 能跨多个网页维持任务，但录像没有给出重复试验次数、失败分布或所有后台调用。

\lecturefigure{slide-32-flight-demo-start.jpg}{订机票 demo 的起始状态}{Stanford 课堂视频 00:23:30}

\begin{knowledgebox}{读图：先确认任务边界而不是先看动画速度}
起始画面只证明演示准备进入航班任务。读者应记录起点账号、目标城市、日期、航空公司偏好和支付权限是否预先配置；这些 hidden state 会显著影响任务难度，也决定结果能否复现。
\end{knowledgebox}

\lecturefigure{slide-33-flight-demo-code.jpg}{通过开发界面向 Agent 发出订票请求}{Stanford 课堂视频 00:23:36}

\begin{knowledgebox}{读图：指令必须转成可执行约束}
自然语言请求通常包含硬约束与软偏好。系统应把日期、机场、人数和预算标成 hard constraints，把 United、出发时段或座位偏好标成 soft preferences；若二者冲突，不能静默牺牲硬约束。
\end{knowledgebox}

\lecturefigure{slide-34-flight-demo-results.jpg}{Agent 浏览航班搜索结果}{Stanford 课堂视频 00:24:33}

\begin{knowledgebox}{读图：搜索结果页是选择问题而非完成状态}
这一帧表明环境已经返回候选航班。Agent 仍需比较价格、转机、时区、行李规则和偏好，并保存“为何选择某项”的 audit trail。页面出现不等于候选已满足用户目标。
\end{knowledgebox}

\lecturefigure{slide-35-flight-demo-checkout.jpg}{Agent 进入 checkout 阶段}{Stanford 课堂视频 00:24:51}

\begin{warningbox}{Checkout 是权限边界}
搜索与填写表单通常可撤销，最终购买可能产生费用和取消成本。进入 checkout 后应展示结构化摘要，并在执行 purchase 前要求显式确认；预存信用卡不能被解释为对任意交易的长期授权。
\end{warningbox}

\lecturefigure{slide-36-flight-demo-itinerary.jpg}{课堂录像中的最终 itinerary 状态}{Stanford 课堂视频 00:25:15}

\begin{knowledgebox}{读图：最终页面仍需要 postcondition 检查}
系统应核对乘客、日期、机场、航班号、总价与 confirmation number，并把结果写入可查询记录。只有 UI 显示 itinerary 但没有核对关键字段，仍可能出现“流程完成、任务错误”的假成功。
\end{knowledgebox}

\teachervoice{讲者称 demo 已个性化到他的 United 偏好、账号和支付方式，并强调 agent 有 purchasing power。课堂最有价值的信号不是炫技，而是它直接把 personalization、credential access 与 irreversible action 绑定在一起。}

\begin{lstlisting}[caption={带确认门的订票状态机}]
state = collect_constraints(request, user_profile)
candidates = search_flights(state)
choice = rank_and_explain(candidates, state)
draft = fill_checkout(choice)
assert validate(draft, state.hard_constraints)
approval = ask_user_confirmation(draft.summary)
if approval:
    receipt = execute_purchase(draft)
    verify_postconditions(receipt, state)
else:
    stop_without_charge()
\end{lstlisting}

\subsection{移动端 Demo：环境约束会中断计划}

第二组画面把同一自然语言控制思路放到手机。platform 首页、消息请求、web search、food result 与 delivery-area failure 组成一条更有教学价值的轨迹，因为最后一帧明确显示环境限制，而不是剪成完全成功的宣传片。本节因此把“失败如何暴露环境约束”作为主线，而不只统计跨了多少应用。

\lecturefigure{slide-37-mobile-platform.jpg}{MultiOn 移动端平台入口}{Stanford 课堂视频 00:25:21}

\begin{knowledgebox}{读图：移动端是新的 I/O 表面}
Agent 要处理小屏 UI、系统权限、通知、app 切换和网络状态。自然语言入口统一了意图表达，却没有统一底层应用状态；每个 app 仍可能拥有不同登录、地理位置和确认规则。
\end{knowledgebox}

\lecturefigure{slide-38-mobile-message.jpg}{通过消息式界面表达任务}{Stanford 课堂视频 00:25:39}

\begin{knowledgebox}{读图：消息是目标，不是完整规范}
一条简短请求常省略预算、份量、地址和替代选项。系统要么从可信 profile 中读取，要么追问；不能用语言模型猜测填补会产生费用或隐私后果的字段。
\end{knowledgebox}

\lecturefigure{slide-39-mobile-web-search.jpg}{移动端 Agent 发起 web search}{Stanford 课堂视频 00:26:03}

\begin{knowledgebox}{读图：跨应用执行需要共享任务状态}
从消息跳到搜索时，目标、约束和来源必须保持一致。若各工具只接收自由文本，信息在多次改写中容易漂移；结构化 task state 能让搜索、筛选与下单共享同一组字段。
\end{knowledgebox}

\lecturefigure{slide-40-mobile-food-result.jpg}{Agent 找到 food ordering 结果}{Stanford 课堂视频 00:26:21}

\begin{knowledgebox}{读图：候选可见仍不等于可交付}
餐厅或商品页只说明搜索成功。真正 postcondition 包括配送地址在范围内、商品可售、费用可接受和预计时间满足要求。Agent benchmark 应把这些环境约束纳入 success definition。
\end{knowledgebox}

\lecturefigure{slide-41-mobile-delivery-limit.jpg}{配送区域限制导致任务受阻}{Stanford 课堂视频 00:26:30}

\begin{warningbox}{失败必须成为显式状态}
Delivery-area constraint 不是让模型“再想久一点”就能消失的推理错误。正确行为是报告阻塞原因、给出可行替代项并等待用户选择，而不是循环点击、隐瞒失败或擅自更改地址。
\end{warningbox}

\teachervoice{这组 demo 最值得保留的是失败帧：环境告诉 agent 当前计划不可行。它提前预示了后面 plan divergence 的主题——没有真实反馈，语言上连贯的计划会越来越远离可执行世界。}

\subsection{Driving-test Claim 与 human-like interface}

课堂随后展示一个 MP4 占位页，讲者称 agent 可完成在线 driving test。由于录像中没有呈现完整题目、答案、评分和重复实验，讲义把它记录为 capability claim，而不把它当作独立验证结果。下一页再从“数字手、眼和人类式接口”解释为何研究者希望 agent 直接使用现有软件。

\lecturefigure{slide-42-driving-test.jpg}{课堂提及的 online driving-test demo}{Stanford 课堂视频 00:26:36}

\begin{warningbox}{Claim 不能替代 evaluation}
要验证 driving-test 能力，至少需要公开题集或封闭测试协议、题目泄漏检查、随机化、多次运行与分项错误分析。一个视频文件图标无法支持通过率、泛化或安全结论。
\end{warningbox}

\lecturefigure{slide-43-human-like-agents.jpg}{Why human-like AI agents：复用为人设计的数字环境}{Stanford 课堂视频 00:28:27}

\begin{knowledgebox}{读图：human-like 指接口兼容，不等于认知等价}
数字环境已有键盘、鼠标、视觉布局和自然语言文档。让 agent 使用这些接口可以复用现有软件，避免为每个网站先开发 API；代价是视觉歧义、点击脆弱性、隐藏状态和更大的攻击面。它描述的是 I/O compatibility，不是说 agent 像人一样理解世界。
\end{knowledgebox}

\subsection{五级 Autonomy：按风险逐步放权}

Autonomy slide 借用自动驾驶分级直觉，把 agent 从人类全程控制逐步推向高自治。具体图中文字较小，教学重点不应是背诵等级名称，而是每一级改变了谁提出计划、谁执行动作、谁监控以及谁承担恢复责任。

\lecturefigure{slide-44-autonomy-levels.jpg}{课堂展示的五级 Agent Autonomy 设想}{Stanford 课堂视频 00:28:45}

\begin{importantbox}{按行动风险设计 autonomy}
\begin{tabularx}{\textwidth}{L{0.10\textwidth}>{\raggedright\arraybackslash}X >{\raggedright\arraybackslash}X}
\toprule
级别 & 系统角色 & 推荐控制 \\
\midrule
L1 & 只建议，用户执行 & 展示依据与可编辑步骤 \\
L2 & 执行低风险可撤销动作 & 每批动作可预览、可取消 \\
L3 & 处理多步任务，关键点确认 & 金钱、外发消息、账号变更前设 gate \\
L4 & 在限定域内自治 & 预算、域名、时间和权限上限 \\
L5 & 广域高度自治 & 课堂愿景；现实部署仍需外部监督与 kill switch \\
\bottomrule
\end{tabularx}
\end{importantbox}

若动作 $a$ 的潜在损失为 $L(a)$、错误概率为 $q(a)$、可恢复系数为 $r(a)\in[0,1]$，可用
\[
R(a)=q(a)L(a)(1-r(a))
\]
作为粗略风险排序。$R(a)$ 越高，越应缩小权限、增加确认或禁止自动执行；该公式是讲义的工程抽象，不是课堂给出的标准。

\teachervoice{讲者把 autonomy 与 self-driving 的分级过程类比，并承认 full autonomy 仍有大量问题。稳健路线不是先宣称“全自动”，而是把可撤销动作与不可逆动作分层，逐步扩大经过验证的自治域。}

\subsection{API 与 direct interaction 的取舍}

Computer Interaction slide 对比两条路线：通过应用提供的 API 操作，或让 agent 直接使用 keyboard、mouse 与屏幕。API 通常结构化、可验证且权限明确；direct interaction 覆盖面更广，却更依赖 UI 稳定性和视觉定位。

\lecturefigure{slide-45-computer-interaction.jpg}{Agent computer interaction 的两条路线}{Stanford 课堂视频 00:31:06}

\begin{knowledgebox}{术语消化：API 与直接操作}
\begin{tabularx}{\textwidth}{L{0.18\textwidth}X X}
\toprule
路线 & 优点 & 主要风险 \\
\midrule
API interaction & 参数与返回值结构化，易做 schema validation、重试和最小权限 & 覆盖受限，需要服务方支持，API 可能变化 \\
Keyboard/mouse & 可复用几乎所有人类界面，原型部署快 & 坐标脆弱、隐藏弹窗、视觉误识别、prompt injection 与误点击 \\
Hybrid & 优先 API，缺失时退回 UI，并共享 task state & 路由与一致性更复杂，必须标记当前执行通道 \\
\bottomrule
\end{tabularx}
\end{knowledgebox}

\teachervoice{课堂明确给出 trade-off：API 更 safe、controllable，direct interaction 更 general。讲义把它转成工程原则——通用性不是免费收益，接口越自由，环境验证与权限隔离越重要。}

\subsection{Action API：把自然语言变成受控动作}

最后一组五张画面演示 Action API 的 overview、代码调用、聊天界面、控制过程和结果。它试图提供一个自然语言层，使开发者不必为每个 UI 动作手写专用自动化；但 API 名称中的 “Action” 不应遮蔽控制流：请求仍需解析、执行、观察、验证和停止。本节按这条控制链逐帧阅读，避免把统一入口误写成统一可靠性。

\lecturefigure{slide-46-action-api-overview.jpg}{MultiOn Action API 的课堂介绍}{Stanford 课堂视频 00:32:18}

\begin{knowledgebox}{读图：统一入口不等于统一语义}
Action API 可以统一调用形式，但底层网站、应用和权限仍不同。可靠实现需要把自然语言转成 typed intent，并保留目标域、允许动作、最大步数和成功条件。
\end{knowledgebox}

\lecturefigure{slide-47-action-api-code.jpg}{开发者通过代码调用 Action API}{Stanford 课堂视频 00:32:21}

\begin{knowledgebox}{读图：SDK 调用只是控制平面的起点}
代码侧至少应记录 request id、policy、budget 与 callback；执行侧返回 structured observations，而不是只回一段“完成了”。这样调用者才能区分运行中、需要确认、失败和已验证完成。
\end{knowledgebox}

\lecturefigure{slide-48-action-api-chat.jpg}{Action API 与聊天式任务输入}{Stanford 课堂视频 00:32:33}

\begin{knowledgebox}{读图：聊天层负责澄清歧义}
当请求缺少对象、时间、预算或收件人时，chat interface 应主动询问。把所有自然语言都直接下发给执行器，会把 conversational ambiguity 转成真实世界错误。
\end{knowledgebox}

\lecturefigure{slide-49-action-api-control.jpg}{Agent 接管应用界面执行动作}{Stanford 课堂视频 00:32:42}

\begin{warningbox}{控制界面时要防止观测与动作错位}
页面滚动、弹窗或异步加载会让旧坐标失效。每个动作前应重新观察目标元素，动作后检查预期状态变化；若 observation 不匹配，应停止而不是沿用过期计划。
\end{warningbox}

\lecturefigure{slide-50-action-api-result.jpg}{Action API demo 的结果状态}{Stanford 课堂视频 00:33:06}

\begin{knowledgebox}{读图：Result 必须绑定 success predicate}
结果画面只有在满足预先声明的 predicate 时才算成功，例如“目标消息已发送给正确收件人且内容一致”。视觉上出现目标 app 或文字，并不足以证明后端状态已经提交。
\end{knowledgebox}

\begin{lstlisting}[caption={Action API 的 observe-act-verify 控制器}]
for step in range(max_steps):
    observation = observe_environment()
    decision = policy.plan(task, observation, permissions)
    if decision.requires_confirmation:
        decision = confirm_with_user(decision)
    result = execute(decision.action)
    if satisfies_success_predicate(result, task):
        return verified_success(result)
    if result.is_blocked or result.risk_exceeded:
        return safe_stop(result)
return timeout_without_claiming_success()
\end{lstlisting}

\subsection{本章小结}

Live demo 能展示 end-to-end 可能性，却不能单独给出普遍成功率。订票流程揭示不可逆购买，移动端失败揭示环境约束，autonomy 分级揭示放权必须按风险推进，Action API 则说明统一自然语言入口仍需要 typed state、confirmation gate 和 postcondition verification。

